\documentclass[11pt]{amsart}
\usepackage{calc,amssymb,amsmath,ulem,stmaryrd,fullpage}
\usepackage{enumerate}
\usepackage{alltt}
\RequirePackage[dvipsnames,usenames]{color}
\let\mffont=\sf
\normalem
%\input{kmacros3.sty}
\input{xy}
\xyoption{all}
\usepackage{hyperref}
\usepackage{graphicx}
\usepackage[all,cmtip]{xy}
\usepackage{verbatim}
\usepackage[left=0.95in,top=0.95in,right=0.95in,bottom=0.95in]{geometry}
\newcommand{\tauCohomology}{T}
\newcommand{\FCohomology}{S}


\theoremstyle{theorem}
%\newtheorem*{mainthm*}{Main Theorem}

\newcommand{\note}[1]{\marginpar{\sffamily\tiny #1}}

\def\todo#1{\textcolor{Mahogany}%
{\sffamily\footnotesize\newline{\color{Mahogany}\fbox{\parbox{\textwidth-15pt}{#1}}}\newline}}

\renewcommand{\O}{\mathcal O}
\newcommand{\ie}{{\itshape i.e.} }
\newcommand{\roundup}[1]{\lceil #1 \rceil}
\newcommand{\rounddown}[1]{\lfloor #1 \rfloor}
\renewcommand{\to}[1][]{\xrightarrow{\ #1\ }}
\newcommand{\onto}[1][]{\protect{\xrightarrow{\ #1\ }\hspace{-0.8em}\rightarrow}}
\newcommand{\into}[1][]{\lhook \joinrel \xrightarrow{\ #1\ }}
%\newcommand{DBDual}[1]{{\underline{\omega}_{#1}}}
\newcommand{\DBDual}[1]{{{\uuline \omega}{}^{\mydot}_{#1}}}
\newcommand{\Tt}{{\mathfrak{T}}}
%\newcommand{\bn}{{\mathfrak{n}} }
\newcommand{\sol}[1]{ \vskip 6pt{\color{blue} {\bf Solution:} #1} \vskip 6pt}

\newcommand{\bR}{{\mathbb{R}}}
\newcommand{\va}{{\vec{a}}}
\newcommand{\vb}{{\vec{b}}}
\newcommand{\vzero}{{\vec{0}}}
\newcommand{\vi}{{\vec{i}}}
\newcommand{\vj}{{\vec{j}}}
\newcommand{\vk}{{\vec{k}}}
\newcommand{\vh}{{\vec{h}}}
\newcommand{\vr}{{\vec{r}}}
\newcommand{\vu}{{\vec{u}}}
\newcommand{\vv}{{\vec{v}}}
\newcommand{\vw}{{\vec{w}}}
\newcommand{\vx}{{\vec{x}}}
\newcommand{\vy}{{\vec{y}}}
\newcommand{\vz}{{\vec{z}}}
\newcommand{\vT}{{\vec{T}}}
\newcommand{\vN}{{\vec{N}}}
\newcommand{\vF}{{\vec{F}}}
\newcommand{\vG}{{\vec{G}}}
\newcommand{\bF}{\mathbb{F}}
\newcommand{\bP}{\mathbb{P}}
\newcommand{\bQ}{\mathbb{Q}}
\newcommand{\bZ}{\mathbb{Z}}
\newcommand{\bA}{\mathbb{A}}
\newcommand{\codim}{{\mathrm{codim}}}
\newcommand{\Map}{{\mathrm{Map}}}
\newcommand{\tld}{\widetilde}
\newcommand{\fram}{\mathfrak{m}}


\begin{document}
\title{Worksheet \#6 -- Math 6130\\Fall 2018}
\author{Due Wednesday, October 31st}

\maketitle
You may work in groups of up to 3.  Only one worksheet needs to be turned in per group.
\vskip 9pt
\noindent
{\bf 1.}  Suppose that $X$ is an affine variety with $R = k[X]$.  Consider an ideal $I \subseteq R$.  Show that the blowup of $I$ is isomorphic to the blowup of $I^2$.
\vskip 3pt
\emph{Hint:} The blowup lives in $X \times \bP^n$ where $I = (f_0, \dots, f_n)$.  $I^2$ is generated by $n^2 +2n + 1$ elements $(f_0^2, f_0 \cdot f_1, \dots, f_n^2)$.  Note on the other hand we have the Segre embedding $\bP^n \to \bP^n \times \bP^n \subseteq P^{N}$ where $N = n^2 + 2n$.
\vskip 220pt
\noindent
{\bf 2.}  Consider $X = \bA^2$ with $k[X] = k[x,y]$.  Show that there is a map from the blowup of $(x^2, xy, y^2)$ (which is isomorphic to the blowup of $(x,y)$ by {\bf 1.}), to the blowup of $(x^2, y^2)$.  Is it an isomorphism?
\newpage
\noindent
{\bf 3.}  Suppose that $X$ is an affine variety and $I, J \subseteq R = k[X]$ are ideals.  Prove that there is a map from $Z$, the blowup of $I \cdot J$,  to $Y$, the blowup of $I$.  $\nu : Z \to Y$.
\vskip 280pt
\noindent
{\bf 4.}  With notation as in {\bf 3.}, after picking generators $f_i$ for $I$ with corresponding affine charts $U_i \subseteq Y$, show $\nu^{-1} (U_i) \to U_i$ is the blowup of the extended ideal $J \cdot k[U_i]$ (extended via the map $k[X] \to k[U_i]$).
\newpage
\noindent
{\bf 5.}  Consider $X = Z(x\cdot y - z^2) \subseteq \bA^3$.  Consider $\pi : Y \to X$ the blowup of the ideal $I = (x,z)$.  Find the locus on which $\pi$ is an isomorphism and show it is strictly larger than $X \setminus Z_X(I)$.
\vskip 250pt
\noindent
{\bf 6.}  Suppose that $X \subseteq Y$ is a closed subvariety of an affine variety.  Consider an ideal $I \subseteq k[x_1, \dots, x_n] = k[\bA^n]$ and let $\pi : \tld{Y} \to Y$ be the blowup of $I$.  Let $\tld{X}$ be the strict transform of $X$ in $Y$.  Prove (without looking \emph{too} closely at the book, which also proves this) that $\pi|_{\tld X} : \tld{X} \to X$ is the blowup of $I \cdot k[X]$.
\newpage
\noindent
{\bf 7.}  Suppose that $I$ is a homogenous prime ideal in $k[x_0, \dots, x_n] = R$ not containing $\fram = (x_0, \dots, x_n)$.  Consider the associated variety $X \subseteq \bA^{n+1}$ (this is called the affine cone over the projective variety $Z = Z_{\bP^n}(I) \subseteq \bP^n$ -- try drawing a picture, because I obviously can't).  Consider the blowup of $\fram$, $\pi : \tld Y \to \bA^{n+1}$ and let $E \cong \bP^n = \pi^{-1} (\text{origin})$ be the exceptional set\footnote{Locus on $\tld Y$ where $\pi$ is not an isomorphism}.  Further let $\tld X \subseteq \tld Y$ be the strict transform of $tld X$.
\vskip 9pt
{\bf{(a)}}  If $I = (f)$ is a principal ideal, consider the affine charts $U_0, \dots, U_n \subseteq \tld Y$ corresponding the generators $x_i$ of $\fram$.  Show that $I \cdot k[U_i]$ can be written as $I_{E \cap U_i}^d \cdot I_{\tld X \cap U_i}$ where $d$ is the degree of $f$ and the ideals $I_{\bullet}$ are taken inside $k[U_i]$.
\vskip 160pt
{\bf{(b)}}  In the setting of (a), show that $U_i \cap E$ are identified with the standard affine charts on $E \cong \bP^n$.  Further show that $Z \cong \tld X \cap E$ and that ${\tld X \cap E \cap U_i}$ correspond to the standard affine charts on $Z$.
\vskip 160pt
{\bf{(c)}}  Show in general that $Z \cong \tld X \cap E$.
\end{document}

